\subsection{Flatness}

\begin{exer}
	Let $M$ be an $A$-module and let $I \subseteq \Ann(m)$ be an ideal.
	\begin{enumerate}[label=(\alph*)]
		\item There is a natural $A/I$-module structure on $M$, and
		$M \cong M \otimes_A A/I$.
		\item Let $N$ be another $A$-module such that
		$I \subseteq \Ann(N)$. The canonical homomorphism
		$M \otimes_A N \to M \otimes_{A/I} N$ is an isomorphism.
	\end{enumerate}	
\end{exer}
\begin{proof}
	For (a), observe that if $x \in I$, then $am = (a + x)m$ for any
	$a \in A$ and $m \in M$. This implies that the multiplication of
	elements of $M$ by cosets $a + I$ is well defined, hence we have an
	$A/I$-module structure on $M$. To see that this is the extension of
	scalars, recall that $M \otimes_A A/I \cong M/(IM),$ which is just $M$.
	\\

	Part $(b)$ follows from a chain of isomorphisms
	\begin{align*}
		& M \otimes_A N\\
		\cong & (M \otimes_A A/I) \otimes_{A/I} N\\
		\cong & M \otimes_{A/I} N,
	\end{align*}
	where the first isomorphism comes from Corollary 1.11 and the second
	comes from part (a).
\end{proof}

\begin{exer}
	Let $\rho: A \to B$ be a ring homomorphism, $S$ a multiplicative subset
	of $A$, and $T = \rho(S)$. Then $T$ is a multiplicative subset of $B$,
	and $T^{-1}B \cong B \otimes_A S^{-1}A$.
\end{exer}
\begin{proof}
	The first part is obvious. To see the last statement, let
	$\iota: B \to T^{-1}B$ be the canonical embedding. Since $\iota$ sends
	elements of $T = \rho(S)$ to units, there is a canonical homomorphism
	$\phi: S^{-1}A \to T^{-1}B$. Since $\otimes_A$ is a pushout in the
	category of $A$-algebras, there is a unique homomorphism
	$B \otimes_A S^{-1}A \to T^{-1}B$ such that this diagram commutes:
	\[
	\begin{tikzcd}
		S^{-1}A \otimes B \arrow[dr, dotted] & S^{-1}A \arrow[l] 
			\arrow[d]\\
		B \arrow[u] \arrow[r] & T^{-1}B.
	\end{tikzcd}
	\]
	This is an isomorphism because of universal property of localizations.
\end{proof}

\begin{exer}
	Nakayama's lemma does not hold for modules that are not finitely
	generated.
\end{exer}
\begin{proof}
	Consider $\QQ$ as a module over $\ZZ_{(p)}$.
\end{proof}

\begin{exer}
	Let $I$ be an ideal of a ring $A$. Then the following are equivalent:
	\begin{enumerate}[label=(\alph*)]
		\item $A/I$ is flat over $A$.
		\item $I^2 = I$
		\item There is an element $e \in I$ such that $e^2 = e$ and
			$eA = I$.
	\end{enumerate}
\end{exer}
\begin{proof}
	\begin{description}[leftmargin=*,font=\normalfont]
		\item[((i)$\Rightarrow$(ii))]
			Consider the exact sequence $0 \to I \to A \to A/I \to
			0$. Since $A/I$ is flat, we have another exact sequence
			$0 \to A/I \to A/I^2 \to 0$, so $I = I^2$.
		\item[((ii)$\Rightarrow$(iii))]
			I don't know how to do this morally. The easy way out
			is to use the form of Nakayama's Lemma proven by the
			determinant trick.
		\item[((iii)$\Rightarrow$(i))]
			Observe that in any maximal ideal, $I$ either localizes
			to $(0)$ or $(1)$. In any case, it's flat at every
			maximal ideal, hence flat overall.
	\end{description}
\end{proof}
